\section{Completion over a fixed seed ball}
\label{sec:completion}

This section supplies necessary conditions for a dephased Hadamard matrix
whose nine seed phases lie in one prescribed ball.  It reduces the remaining
completion problem to a covering of $U(2)$.  The reduction uses frames fixed
at the seed center; no exact Hadamard matrix is assumed to occur at that
center.  All statements in this section first concern exact real and
complex arithmetic.  The implementation uses the outward replacements and
comparison margins justified in Section~\ref{sec:rounding}.

\subsection{Seed blocks and fixed complementary frames}

Fix the stored seed center, taking its binary phase coordinates as exact
real numbers.  After moving the shared column into the second position,
write its relevant unit entries as $a,b_0,\ldots,b_3,c_0,\ldots,c_3$, and put
\begin{equation}
\label{completion:seed-blocks}
 E_0=\begin{pmatrix}1&1\\1&a\end{pmatrix},\qquad
 B_0=\begin{pmatrix}1&1&1&1\\b_0&b_1&b_2&b_3\end{pmatrix},\qquad
 (C_0)_{i,:}=(1,c_i).
\end{equation}
A feasible completion is an entrywise unit-modulus matrix
\[
 H'=\begin{pmatrix}E'&B'\\ C'&D'\end{pmatrix},\qquad
 H'H'^*=H'^*H'=6I,
\]
with first row and first column equal to one.  Its seed displacements,
measured in radians, are ordered as
\[
 s=(s_a,s_{b_0},\ldots,s_{b_3},s_{c_0},\ldots,s_{c_3})\in\R^9.
\]
Thus the varied entries are $a e^{i s_a}$, $b_j e^{i s_{b_j}}$, and
$c_i e^{i s_{c_i}}$.  The shared entry occurs once in $s$.  With zero-based
row-major phase indexing, these coordinates occupy
\begin{equation}
\label{completion:position-map}
 (7,8,9,10,11,13,19,25,31).
\end{equation}
The enlarged seed radius in cycles is denoted by $r_s$, and
$R_s=2\pi r_s$.  In the stated computation $r_s=.0320001$.
Set $\Delta E=E'-E_0$, and define $\Delta B,\Delta C$ similarly.  To avoid
confusing norms with entries of the blocks, write
\[
 e_1=\lVert\Delta E\rVert_F,\qquad
 b_1'=\lVert\Delta B\rVert_F,\qquad
 c_1'=\lVert\Delta C\rVert_F.
\]
The chord bound $|e^{ix}-1|\le |x|$ gives
\begin{equation}
\label{completion:seed-norm}
 e_1^2+(b_1')^2+(c_1')^2\le\lVert s\rVert^2\le R_s^2.
\end{equation}

\begin{lemma}[Fixed frames]
\label{completion:frames}
Let $p\in\C^4$ have unit-modulus entries, let
$S=\sum_j p_j$, and set $M=(\mathbf 1\ p)$.  The squared singular values
of $M$ are $4\pm |S|$.  If $|S|<4$, then
\begin{equation}
\label{completion:inverse-frame}
 M(M^*M)^{-1}
 =\frac1{16-|S|^2}
 \big(4\mathbf 1-p\overline S\,,\;4p-\mathbf 1 S\big).
\end{equation}
Under the same condition $|S|<4$, there is an orthonormal $4$-by-$2$
matrix $K$ with range $\ker M^*$.
It may be constructed by completing $\mathbf 1/2$ and the normalized
residual of $p$ by projected coordinate vectors with nonzero residuals.
\end{lemma}

\begin{proof}
Direct multiplication gives
\[
 M^*M=\begin{pmatrix}4&S\\\overline S&4\end{pmatrix}.
\]
Its eigenvalues and inverse give the assertions about singular values and
\eqref{completion:inverse-frame}.  The first two Gram--Schmidt vectors
span $\operatorname{ran}M$.  At each subsequent step, the orthogonal
complement of the vectors already selected is nonzero.  Since the coordinate
vectors span $\C^4$, at least one has a nonzero projection into that
complement.  Selecting and normalizing such a projection gives the next
orthonormal vector.  The final two vectors form $K$.
\end{proof}

Apply Lemma~\ref{completion:frames} to $p=(\overline b_j)_j$ and
$p=(c_i)_i$.  Throughout the seed ball keep the resulting frames fixed, and
write
\begin{equation}
\label{completion:frame-definitions}
 G_B=B_0^*(B_0B_0^*)^{-1},\qquad
 G_C=C_0(C_0^*C_0)^{-1},\qquad
 \operatorname{ran}K_B=\ker B_0,\quad
 \operatorname{ran}K_C=\ker C_0^*.
\end{equation}
Here $K_B^*K_B=K_C^*K_C=I_2$.  Define
\[
 m_B=\sigma_{\min}(B_0),\quad
 m_C=\sigma_{\min}(C_0),\quad
 m=\min(m_B,m_C)>0,
\]
so that $\lVert G_B\rVert_{\rm op}=1/m_B$ and
$\lVert G_C\rVert_{\rm op}=1/m_C$.
The numerical construction chooses particular coordinate vectors; the ideal
frames in this argument use those same choices.  The acceptance checks on
rank and residual norms certify that this construction is defined.  A
failed frame check leaves the seed unresolved and never excludes it.

We shall use the orthogonal projections
\begin{equation}
\label{completion:projections}
 P_K=K_BK_B^*,\qquad P_B=I-P_K=G_BB_0,\qquad
 P_C=I-K_CK_C^*=G_CC_0^*.
\end{equation}
Finally, put
\begin{equation}
\label{completion:A-bounds}
 A_0=\lVert E_0\rVert_{\rm op}=\sqrt{2+|1+a|},\qquad
 A_* =\min(2,A_0+R_s).
\end{equation}
Both $A_0$ and $\lVert E'\rVert_{\rm op}$ are at most $A_*$.
Indeed $A_0\le2$, every entry of $E'$ has modulus one and hence
$\lVert E'\rVert_{\rm op}\le\lVert E'\rVert_F=2$, and
$\lVert E'\rVert_{\rm op}\le A_0+e_1\le A_0+R_s$.

\subsection{A unitary associated to every feasible completion}

The polar decomposition associates a unitary with a nonsingular square
matrix, and its nearest-unitary property in unitarily invariant norms is
classical~\cite[Theorem~1]{FanHoffman1955}; see also
\cite[Theorem~1.1 and Corollary~2.3]{Higham1986Polar}. Here the compressed
completion has an explicit positive defect. The next lemma derives the
trace-defect estimate required for the finite $U(2)$ cover.


\begin{lemma}[Polar comparison]
\label{completion:polar}
For a feasible completion let $Y=K_C^*D'K_B$.  Then
\begin{align}
\label{completion:deficit-identity}
 6I-Y^*Y
 &=K_B^*B'^*B'K_B+
   K_B^*D'^*P_CD'K_B\succeq0,\\
\label{completion:deficit-bound}
 \operatorname{tr}(6I-Y^*Y)
 &\le (b_1')^2+
 \frac{(\lVert E'\rVert_{\rm op}b_1'+\sqrt6\,c_1')^2}{m_C^2}
 \le q,
\end{align}
where
\begin{equation}
\label{completion:q-and-p}
 q=R_s^2\left(1+\frac{A_*^2+6}{m_C^2}\right).
\end{equation}
If $q<6$, the polar factor
$V=Y(Y^*Y)^{-1/2}$ belongs to $U(2)$ and satisfies
\begin{equation}
\label{completion:polar-bound}
 \lVert Y-\sqrt6 V\rVert_F\le
 p:=\frac{q}{\sqrt6+\sqrt{6-q}}.
\end{equation}
\end{lemma}

\begin{proof}
The lower-right block of $H'^*H'=6I$ is
$B'^*B'+D'^*D'=6I$.  Compress it by $K_B$, and subtract
$Y^*Y=K_B^*D'^*(I-P_C)D'K_B$.  This proves
\eqref{completion:deficit-identity}.  Its trace is
\[
 \lVert B'K_B\rVert_F^2+\lVert P_CD'K_B\rVert_F^2.
\]
Since $B_0K_B=0$, the first term is at most $(b_1')^2$.
The off-diagonal block of column orthogonality gives
$C'^*D'=-E'^*B'$, and consequently
\[
 C_0^*D'K_B=-E'^*\Delta B K_B-\Delta C^*D'K_B.
\]
Now $P_C=G_CC_0^*$ and $\lVert D'\rVert_{\rm op}\le\sqrt6$.
These facts give the first inequality in \eqref{completion:deficit-bound}.
For the second, apply Cauchy--Schwarz to
$A_*b_1'+\sqrt6 c_1'$ and then use
\eqref{completion:seed-norm}.

Let $\ell_1,\ell_2$ be the eigenvalues of $6I-Y^*Y$.
They are nonnegative and have sum at most $q<6$.  The singular values of
$Y$ are therefore positive and equal to $\sqrt{6-\ell_j}$, so its polar
factor is unitary.  Moreover
\[
 \sqrt6-\sqrt{6-\ell_j}
 =\frac{\ell_j}{\sqrt6+\sqrt{6-\ell_j}}
 \le\frac{\ell_j}{\sqrt6+\sqrt{6-q}}.
\]
Taking the Euclidean norm of these two singular-value differences, and
using $\sqrt{\ell_1^2+\ell_2^2}\le\ell_1+\ell_2\le q$, proves
\eqref{completion:polar-bound}.
\end{proof}

The inequality $q<6$ is checked using outward bounds.  In the specified
parameter range the stronger bound $q<.32$ follows from the accepted frame
checks and the seed radius; see Section~\ref{sec:rounding}.  The reduction
therefore assigns at least one $V\in U(2)$ to every feasible completion of
an accepted seed ball.

\subsection{The unitary cells and their Taylor bounds}

\begin{lemma}[A closed covering of $U(2)$]
\label{completion:unitary-cover}
Every $V\in U(2)$ has a representation
\begin{equation}
\label{completion:unitary-parameterization}
 V=e(\alpha)
 \begin{pmatrix}
 e(\beta)\cos(2\pi t)&e(\gamma)\sin(2\pi t)\\
 -e(-\gamma)\sin(2\pi t)&e(-\beta)\cos(2\pi t)
 \end{pmatrix},
\end{equation}
with
$0\le\alpha\le\tfrac12$, $0\le t\le\tfrac14$, and
$0\le\beta,\gamma\le1$.
For an integer $N_u$ divisible by four, the products of closed intervals
of width $1/N_u$ in these ranges cover this parameter domain.  Their four
bin counts are $N_u/2,N_u/4,N_u,N_u$.
\end{lemma}

\begin{proof}
Choose a square root of $\det V$ having phase $\alpha\in[0,1/2]$ in
cycles.  Then $e(-\alpha)V\in SU(2)$ and has the form
\[
 \begin{pmatrix}z&w\\-\overline w&\overline z\end{pmatrix},
 \qquad |z|^2+|w|^2=1.
\]
Choose $t\in[0,1/4]$ with
$|z|=\cos(2\pi t)$ and $|w|=\sin(2\pi t)$, and choose phases $\beta$
and $\gamma$ for $z$ and $w$.  A phase of a zero entry may be chosen
arbitrarily.  Thus no boundary case requires division by an entry.
The interval products cover the stated closed domain directly.
\end{proof}

Fix one cell and denote its midpoint matrix by $U$.  If $V$ is represented
inside the cell, let
\begin{equation}
\label{completion:unitary-cell}
 \theta=(\theta_\alpha,\theta_t,\theta_\beta,\theta_\gamma),\qquad
 |\theta_j|\le h:=\frac{\pi}{N_u},
\end{equation}
be its displacement from the midpoint in radians.  Denote the midpoint
angular differential, applied to this displacement, by $dU[\theta]$.
All four derivatives here are with respect to radian angles, although
\eqref{completion:unitary-parameterization} uses cycle coordinates.

\begin{lemma}[Uniform cell Taylor estimates]
\label{completion:unitary-taylor}
For the cell and displacement in \eqref{completion:unitary-cell},
\begin{equation}
\label{completion:unitary-taylor-bounds}
 \lVert V-U\rVert_F\le\sqrt6 h,\qquad
 \lVert V-U-dU[\theta]\rVert_F\le\frac{\sqrt{42}}2h^2.
\end{equation}
\end{lemma}

\begin{proof}
Follow the straight segment in the four angle coordinates and write its
matrix as $U(\tau)$, $0\le\tau\le1$.  For this proof abbreviate the four
constant angular velocities as
$x=\theta_\alpha$, $y=\theta_t$, $z=\theta_\beta$, and
$w=\theta_\gamma$, and let $T$ be the current radian $t$ angle.
Differentiating the four entries yields
\[
 \lVert U'(\tau)\rVert_F^2
 =2\big(x^2+y^2+z^2\cos^2 T+w^2\sin^2 T\big)\le6h^2.
\]
For the two diagonal entries, the squared magnitudes of the real factors
in the second derivative have sum
\begin{align*}
 &\big(((x+z)^2+y^2)^2+((x-z)^2+y^2)^2\big)\cos^2 T\\
 &\qquad=2\big((x^2+z^2+y^2)^2+4x^2z^2\big)\cos^2 T
 \le26h^4\cos^2 T.
\end{align*}
The imaginary factors of those two entries contribute
\[
 4y^2\big((x+z)^2+(x-z)^2\big)\sin^2 T
 \le16h^4\sin^2 T.
\]
The off-diagonal entries give the same two bounds with $z$ replaced by $w$
and cosine and sine interchanged.  Real and imaginary factors add in
squared magnitude.  Hence $\lVert U''(\tau)\rVert_F^2\le42h^4$.
Integrating the first derivative, and then using the integral Taylor
remainder with weight $1-\tau$, proves
\eqref{completion:unitary-taylor-bounds}.
\end{proof}

For clarity about the implementation, if the entries of $U$ are $U_{ij}$,
its alpha derivative is $iU$, its beta derivative has diagonal entries
$iU_{00},-iU_{11}$ and zero off-diagonal entries, and its gamma derivative
has off-diagonal entries $iU_{01},-iU_{10}$ and zero diagonal entries.
The $t$ derivative differentiates the displayed sine and cosine factors.
These are precisely the four angular derivatives assembled by
\texttt{uc::init} and \texttt{uc::initDeriv}.

\subsection{Comparison with a cell center}

For the fixed seed and unitary cell define
\begin{equation}
\label{completion:comparison-def}
 D=D_{\rm span}+D_{\rm ker},\qquad
 D_{\rm span}=-C_0E_0^*G_B^*,\qquad
 D_{\rm ker}=\sqrt6 K_C U K_B^*,\qquad Z=D'-D.
\end{equation}
The symbol $D$ henceforth denotes this comparison block, which need not be
entrywise flat.  It is not asserted to be the lower block of a Hadamard
matrix at the seed center.

\begin{lemma}[Orthogonal comparison pieces]
\label{completion:comparison-pieces}
The matrices $ZP_B$, $P_CZP_K$, and $(I-P_C)ZP_K$ are pairwise
Frobenius-orthogonal and sum to $Z$.  They satisfy
\begin{align}
\label{completion:piece-B}
 ZP_B&=(-\Delta C E_0^*-C'\Delta E^*-D'\Delta B^*)G_B^*,\\
\label{completion:piece-C}
 P_CZP_K&=G_C(-E'^*\Delta B-\Delta C^*D')P_K,\\
\label{completion:piece-K}
 (I-P_C)ZP_K&=K_C(Y-\sqrt6 U)K_B^*.
\end{align}
Consequently every feasible completion associated to this unitary cell
satisfies
\begin{equation}
\label{completion:z-bound}
 \lVert Z\rVert_F\le z:=
 \sqrt{\left(\frac{R_s(\sqrt6+A_*)}{m}\right)^2+(6h+p)^2}.
\end{equation}
\end{lemma}

\begin{proof}
The right projections $P_B,P_K$ are complementary and orthogonal; within
the right $P_K$ part, the left projections $P_C,I-P_C$ are complementary
and orthogonal.  This proves the decomposition and orthogonality.
Row orthogonality gives $D'B'^*=-C'E'^*$.  Since
$P_B=B_0^*G_B^*$, $D_{\rm ker}P_B=0$, and
$D_{\rm span}P_B=D_{\rm span}$, subtracting the comparison block gives
\eqref{completion:piece-B}.  The identity
$C'^*D'=-E'^*B'$, together with $B_0P_K=0$ and
$P_CD P_K=0$, gives \eqref{completion:piece-C}.
The definitions of $Y$, $K_B$, and $K_C$ give
\eqref{completion:piece-K}.

Put $r_{EB}=\sqrt{e_1^2+(b_1')^2}$.  The four lower rows of $H'$ are
orthogonal and have norm $\sqrt6$, so
$\lVert[C'\ D']\rVert_{\rm op}=\sqrt6$.  Therefore
\begin{align*}
 \lVert ZP_B\rVert_F
 &\le\frac{\sqrt6 r_{EB}+A_0c_1'}{m_B},\\
 \lVert P_CZP_K\rVert_F
 &\le\frac{A_*b_1'+\sqrt6 c_1'}{m_C}
 \le\frac{A_*r_{EB}+\sqrt6 c_1'}m.
\end{align*}
The vector of these two norms is bounded componentwise by
\[
 \frac1m
 \begin{pmatrix}\sqrt6&A_*\\ A_*&\sqrt6\end{pmatrix}
 \begin{pmatrix}r_{EB}\\c_1'\end{pmatrix}.
\]
The displayed matrix has operator norm $\sqrt6+A_*$; its input norm is at
most $R_s$ by \eqref{completion:seed-norm}.
For the third piece, multiplication by the two orthonormal frames preserves
Frobenius norm, and
\[
 \lVert Y-\sqrt6 U\rVert_F
 \le\lVert Y-\sqrt6 V\rVert_F+\sqrt6\lVert V-U\rVert_F
 \le p+6h.
\]
Pythagoras for the three pieces gives \eqref{completion:z-bound}.
\end{proof}

\begin{lemma}[A row-block operator bound]
\label{completion:L-bound}
For every cell midpoint $U\in U(2)$,
\begin{equation}
\label{completion:L-definition}
 \lVert[C_0\ D]\rVert_{\rm op}\le
 L:=\max\left(\sqrt6,\;
 \lVert C_0\rVert_{\rm op}\sqrt{1+(A_0/m_B)^2}\right).
\end{equation}
Also $\lVert DP_K\rVert_{\rm op}=\sqrt6$ and
\begin{equation}
\label{completion:D-operator-bound}
 \lVert D\rVert_{\rm op}
 \le\max\left(\sqrt6,\frac{\lVert C_0\rVert_{\rm op}A_0}{m_B}\right).
\end{equation}
\end{lemma}

\begin{proof}
The two terms $D_{\rm span},D_{\rm ker}$ in
\eqref{completion:comparison-def} have orthogonal right domains and
orthogonal ranges.  In particular the cross terms in $DD^*$ vanish, and
$G_B^*G_B=(B_0B_0^*)^{-1}$ gives
\begin{equation}
\label{completion:L-Gram}
 [C_0\ D][C_0\ D]^*
 =C_0\big[I+E_0^*(B_0B_0^*)^{-1}E_0\big]C_0^*
   +6K_CK_C^*.
\end{equation}
The first summand has range in $\operatorname{ran}C_0$ and operator norm at
most $\lVert C_0\rVert_{\rm op}^2(1+A_0^2/m_B^2)$.  The second is six
times the projection onto its orthogonal complement.  Thus their sum has
norm equal to the larger of their norms, proving
\eqref{completion:L-definition}.  Omitting the $C_0C_0^*$ contribution
proves \eqref{completion:D-operator-bound}.  Finally $DP_K=D_{\rm ker}$,
whose two nonzero singular values are $\sqrt6$.
\end{proof}

\subsection{A uniform linearization remainder}

Let $E_1,B_1,C_1$ be the seed first variations: at a varied entry, multiply
the center entry by $i$ times the corresponding component of $s$; all fixed
entries have variation zero.  Define the linear map
\begin{equation}
\label{completion:linear-map}
 \begin{split}
 \mathcal J(s,\theta)={}&
 (-C_1E_0^*-C_0E_1^*-DB_1^*)G_B^*\\
 &+G_C(-E_0^*B_1-C_1^*D)P_K
 +\sqrt6 K_CdU[\theta]K_B^*.
 \end{split}
\end{equation}
This is a necessary linear comparison inferred from the exact
orthogonality equations.  In particular it is not obtained by assuming
that the fixed frames vary differentiably with the seed or that an exact
completion exists at every seed center.

\begin{lemma}[Linear remainder bound]
\label{completion:linear-remainder}
Every feasible completion associated to the cell satisfies
\begin{equation}
\label{completion:linear-remainder-bound}
 \lVert Z-\mathcal J(s,\theta)\rVert_F\le E_r,
 \qquad
 E_r=\sqrt{e_{\rm off}^2+
       \left(\frac{\sqrt{6\cdot42}}2h^2+p\right)^2},
\end{equation}
where
\begin{equation}
\label{completion:e-off}
 e_{\rm off}=
 \frac{R_s^2(L+A_0+1)}{2m}+\frac{zR_s}{m}.
\end{equation}
\end{lemma}

\begin{proof}
The first two terms of \eqref{completion:linear-map} lie respectively in
the spaces of the first two projections in
Lemma~\ref{completion:comparison-pieces}; the last term lies in the third.
We bound the remainders in those spaces.

First replace $C',D'$ in \eqref{completion:piece-B} by $C_0,D$.
The additional numerator terms are
$-\Delta C\Delta E^*-Z\Delta B^*$.
Replacing $E',D'$ in \eqref{completion:piece-C} by $E_0,D$ leaves
$-\Delta E^*\Delta B-\Delta C^*Z$.
Across the two orthogonal output spaces, the seed-only product terms have
combined Frobenius norm at most
\[
 \frac{e_1\sqrt{(b_1')^2+(c_1')^2}}m
 \le\frac{R_s^2}{2m}.
\]
The product terms involving $Z$ have combined norm at most
\[
 \frac{\lVert Z\rVert_F\sqrt{(b_1')^2+(c_1')^2}}m
 \le\frac{zR_s}{m}.
\]
These inequalities use submultiplicativity with a Frobenius norm on the
varying factor, the bounds on $G_B,G_C$, and contraction by $P_K$.

It remains in the first two spaces to replace each of
$\Delta E,\Delta B,\Delta C$ by its first variation.
Since $|e^{ix}-1-ix|\le x^2/2$, the combined Frobenius norm of their three
Taylor remainders is at most
\[
 \frac12\left(\sum_{j=1}^9s_j^4\right)^{1/2}
 \le\frac12\lVert s\rVert^2\le\frac12R_s^2.
\]
For arbitrary remainder inputs $F_E,F_B,F_C$, put
$x=(\lVert F_E\rVert_F^2+\lVert F_B\rVert_F^2)^{1/2}$ and
$y=\lVert F_C\rVert_F$.  Lemma~\ref{completion:L-bound} shows that the
norms of their two linear output pieces are bounded componentwise by
\[
 \frac1m\begin{pmatrix}L&A_0\\ A_0&\sqrt6\end{pmatrix}
 \begin{pmatrix}x\\y\end{pmatrix}.
\]
For the second component use $\lVert DP_K\rVert_{\rm op}=\sqrt6$.
Since $L\ge\sqrt6$, the matrix has operator norm at most $L+A_0$.
The combined seed Taylor contribution is therefore at most
$R_s^2(L+A_0)/(2m)$.  Adding the two product bounds gives
\eqref{completion:e-off}.

In the third space the exact remainder is
\[
 K_C\big[Y-\sqrt6 U-\sqrt6 dU[\theta]\big]K_B^*.
\]
Insert $\sqrt6 V$, and apply the polar and unitary Taylor bounds.  Its
norm is at most $p+\sqrt{6\cdot42}h^2/2$.
This space is orthogonal to the first two, so Pythagoras proves
\eqref{completion:linear-remainder-bound}.
\end{proof}

Here is the entrywise form used to assemble \eqref{completion:linear-map}.
For $0\le i,j,k<4$, its shared seed coefficient is
\[
 -c_i\overline{i a(G_B)_{j1}},
\]
its $b_k$ coefficient is
\[
 -\big[(G_C)_{i0}+(G_C)_{i1}\overline a\big]\,
       i b_k(P_K)_{kj}
 -D_{ik}\overline{i b_k(G_B)_{j1}},
\]
and its $c_k$ coefficient is
\[
 -\mathbf 1_{i=k}\,i c_k
       \big(\overline{(G_B)_{j0}}+\overline{a(G_B)_{j1}}\big)
 -(G_C)_{i1}\overline{i c_k}(D_{\rm ker})_{kj}.
\]
For the angular coefficients, if the entries of $U$ are flattened in
row-major order, use
\[
 W_{ij,l}=\sqrt6(K_C)_{i,\lfloor l/2\rfloor}
                 \overline{(K_B)_{j,l\bmod2}},\qquad 0\le l<4,
\]
and contract with each angular derivative of $U$.
These formulas identify the fixed and moving parts of
\texttt{block::setup}, \texttt{makeD}, and \texttt{makeJ}.

\subsection{Necessary flatness and separating inequalities}

We have bounded both the total block displacement $Z=D'-D$ by $z$ and
its remainder after the linear comparison $\mathcal J(s,\theta)$ by
$E_r$.  We now use these two bounds in the necessary equations
$|D'_{ij}|=1$: the first test uses the total displacement alone, and the
second also uses its linear comparison and the seed sum equations.

\begin{lemma}[Norm exclusion]
\label{completion:norm-exclusion}
A feasible completion associated to a cell must satisfy
\begin{equation}
\label{completion:norm-test}
 \sum_{i,j=0}^3(|D_{ij}|-1)^2\le z^2.
\end{equation}
\end{lemma}

\begin{proof}
Each $D'_{ij}$ has modulus one, and the reverse triangle inequality gives
$\big||D_{ij}|-1\big|\le|D'_{ij}-D_{ij}|$.
Sum the squares and use \eqref{completion:z-bound}.
\end{proof}

For the second test define two complex constants and their linear parts by
\begin{align*}
 c_X&=(1+a+\textstyle\sum_j b_j)/i,&
 J_Xs&=a s_a+\textstyle\sum_j b_j s_{b_j},\\
 c_Y&=(1+a+\textstyle\sum_i c_i)/i,&
 J_Ys&=a s_a+\textstyle\sum_i c_i s_{c_i}.
\end{align*}
The corresponding exact row and column sums vanish.  Taylor's formula
therefore gives $c_X+J_Xs+\rho_X=0$ and
$c_Y+J_Ys+\rho_Y=0$, with
\begin{equation}
\label{completion:seed-sum-remainders}
 |\rho_X|\le\frac12R_s^2,\qquad
 |\rho_Y|\le\frac12R_s^2.
\end{equation}
Take real and imaginary parts of these two equations, in that order.
Append the following sixteen real constants and linear forms:
\begin{equation}
\label{completion:flat-equations}
 c_{ij}=\frac{|D_{ij}|^2-1}{2},\qquad
 J_{ij}(s,\theta)
 =\operatorname{Re}\big(\overline{D_{ij}}\mathcal J_{ij}(s,\theta)\big).
\end{equation}
Denote the resulting real constant vector by $c_{\rm all}\in\R^{20}$
and the real linear matrix by $J_{\rm all}\in\R^{20\times13}$, with
its first nine columns corresponding to $s$ and its last four to $\theta$.

\begin{lemma}[Directional exclusion]
\label{completion:directional-exclusion}
Let $y\in\R^{20}$ be arbitrary, split as two pairs $y_X,y_Y\in\R^2$
and a flat part $y_D\in\R^{16}$, and write
\[
 g=J_{\rm all}^*y=(g_s,g_\theta)\in\R^9\times\R^4.
\]
Every feasible completion associated to the cell satisfies
\begin{equation}
\label{completion:directional-test}
 \begin{split}
 y\cdot c_{\rm all}\le{}&
 R_s\lVert g_s\rVert+h\lVert g_\theta\rVert_1
 +E_r\lVert D\circ y_D\rVert_F\\
 &+\frac12z^2\max\big(0,\max_{i,j}(-y_{D,ij})\big)
 +\frac12R_s^2(\lVert y_X\rVert+\lVert y_Y\rVert).
 \end{split}
\end{equation}
Here $\circ$ denotes entrywise multiplication, with the sixteen flat
coordinates arranged as a $4$-by-$4$ matrix.
\end{lemma}

\begin{proof}
Write $R=Z-\mathcal J(s,\theta)$, so
$\lVert R\rVert_F\le E_r$.  Expanding the exact flatness equation gives
\[
 0=\frac{|D'_{ij}|^2-1}{2}
   =c_{ij}+J_{ij}(s,\theta)
    +\operatorname{Re}(\overline{D_{ij}}R_{ij})
    +\frac12|Z_{ij}|^2.
\]
Weight these equations by $y_D$, and weight the two complex seed equations
by $y_X,y_Y$ using their real coordinate pairs.  Move every term other than
$y\cdot c_{\rm all}$ to the right.
The linear contribution is bounded above by
$R_s\lVert g_s\rVert+h\lVert g_\theta\rVert_1$, using the seed ball and
angle box.  Real Cauchy--Schwarz bounds the flat remainder contribution by
$E_r\lVert D\circ y_D\rVert_F$.  The seed remainders contribute at most
the last term of \eqref{completion:directional-test} by
\eqref{completion:seed-sum-remainders}.

The quadratic term has a fixed sign before weighting.  After it is moved
to the right its contribution is
\[
 -\frac12\sum_{i,j}y_{D,ij}|Z_{ij}|^2
 \le\frac12\max\big(0,\max_{i,j}(-y_{D,ij})\big)
               \sum_{i,j}|Z_{ij}|^2.
\]
Use $\lVert Z\rVert_F\le z$.  This proves the remaining term and the
stated inequality.
\end{proof}

\begin{corollary}[Dual certificates]
\label{completion:dual-certificates}
A cell contains no feasible completion if
\eqref{completion:norm-test} fails, or if some real direction $y$ violates
\eqref{completion:directional-test}.  The choice of $y$ may be heuristic;
only its final inequality must be certified.  Positive normalization of
$y$ preserves the certificate.
\end{corollary}

\begin{proof}
Both inequalities are necessary for every feasible completion in the cell.
All terms of \eqref{completion:directional-test} are positively homogeneous
in $y$.  Thus a strict violation is a separating, or dual, certificate,
independent of how the direction was found.
\end{proof}

The implementation stacks exactly these twenty equations in \texttt{eq}.
The routine \texttt{trial::pass} proposes directions by regularized linear
solves and weight updates.  It excludes only when
\texttt{trial::certificate} verifies the final inequality.  That routine
checks finiteness, divides by the largest absolute coordinate, and treats
the resulting stored coefficients as the direction.  It uses the negative
part of the flat weights, as required by
\eqref{completion:directional-test}.  No accuracy or convergence assertion
about the proposing solves is needed.  The used Stage A path is
\texttt{Asearch::testA}; the auxiliary \texttt{Usearch::test} and
\texttt{scalar} routines do not supply additional exclusion tests for this
run.  The actual norm and directional comparisons include the margins
proved in Section~\ref{sec:rounding}.

\subsection{Finite recursion and unresolved cells}

\begin{proposition}[Completion coverage conditional on successful termination]
\label{completion:recursion}
Fix a seed ball for which the frame and bound constructions are certified.
Suppose that every final unitary cell is either excluded by a certified
necessary-condition test or covered by a certified phase ball as in
Section~\ref{sec:phase-cover}.  Then those phase balls cover all feasible
Hadamard completions of the seed ball, with the equivalences specified by
the phase-ball certificates.  The Stage A recursion establishes precisely
this conclusion if it terminates with zero unresolved cells.
\end{proposition}

\begin{proof}
Given any feasible completion, Lemma~\ref{completion:polar} produces a
unitary $V$.  Lemma~\ref{completion:unitary-cover} places one of its
parameter representations in a closed base cell.  Every certified
exclusion test is necessary for that completion, so it cannot exclude a
cell containing this representation.  If a cell is subdivided, the two
closed children in each coordinate cover the parent interval; their
sixteen products cover the parent cell.  Hence the representation can be
followed into at least one child.  A successful phase-ball certificate
covers every feasible completion associated to its entire cell, so it
covers the chosen completion if that outcome is reached.

This is a finite induction through the processed cell tree.  Zero
unresolved cells means that every branch ends in a certified exclusion or
a certified phase-ball assignment.  The branch followed by a feasible
completion cannot end in exclusion, and therefore ends in coverage.
\end{proof}

In the stated invocation the base grid is $N_u=16$.  Retained cells are
split to $N_u=32$ and then $64$ by replacing every bin index $b_j$ with
both $2b_j$ and $2b_j+1$.  At the first level only the norm test is needed.
At $64$ and deeper, every unexcluded cell is offered to the phase-ball
coverage routine; unsuccessful coverage causes all sixteen children to be
processed with the same seed and the same fixed frames, using the new
value of $h$.  Subdivision is allowed through $N_u=512$.

A coverage or snapping failure does not exclude a cell.  A live cell at the
depth cap increments the unresolved count.  Exceeding the deeper-node cap
of $1{,}000{,}000$ sets the give-up flag and increments that count; later
short returns under that flag cannot restore successful status.  A failed
frame construction is recorded as a failed seed.  The driver records every
seed with unresolved cells and returns failure if any seed failed.  Thus a
negative or global-coverage conclusion requires successful processing of
all prescribed seed indices, including all shards, with no failed seed or
unresolved cell.  Retaining some valid phase balls from an unsuccessful
seed does not certify the unprocessed remainder.  These conditions are
part of the computational obligation, not consequences of the local
inequalities alone.
